\section{Structure below the provision}
\label{sec:structure}

Cross-references in dense documents rarely point at a whole provision. They
point inside one: ``Article 66(2)(e)(iv)'', ``paragraph 1, point (b)'', ``the
second subparagraph of paragraph 3''. Readers ask questions the same way. To
resolve either, every passage has to know not only its article but its
paragraph, point, sub-point and subparagraph. This was the first thing it took
to make the system output cross-references: until passages carry these labels,
a reference can be resolved only to a whole provision, and ``Article 66(2)(e)''
cannot be told apart from Article 66.

\subsection{Paragraphs, points and sub-points}

Inside a provision three kinds of marker open a line: a numbered paragraph
(``2.''), a lettered point (``(b)'' or ``b.'') and a roman sub-point (``(iv)''
or ``iv.''). The shapes alone are ambiguous: ``(i)'', ``(v)'' and ``(x)'' are
both letters and roman numerals, and a number followed by a full stop can begin
a wrapped line in the middle of a sentence.

The parser therefore resolves markers by \emph{sequence}. After each heading it
expects paragraph 1 and point (a); it then tracks the next paragraph number,
the next point letter and, once a point is open, the next roman numeral. A
marker is accepted only if it is the one expected next (\cref{lst:sequence}).
After point (h), ``(i)'' is the expected letter and becomes point (i). After
point (k) of Article 7(2) of the AI Act---``the extent to which existing Union
law provides for:''---the expected letter is (l), so the ``(i)'' that follows
is taken as the first sub-point of (k). A paragraph number is accepted only if
it is the next one expected and the text before it ends a sentence or a clause
(with ``.'', ``;'' or ``:''), which rejects numbers that merely begin a line.

\begin{lstlisting}[float, language=Python, caption={Accepting a point or sub-point
only in sequence (\texttt{ingestion/structure.py}, abridged).},
label={lst:sequence}]
elif kind == "marker" and next_paragraph is not None:
    if value == next_point:
        marks.append((pos, "point", value))
        next_point = chr(ord(value) + 1)
        next_subpoint = "i"
    elif value == next_subpoint:
        marks.append((pos, "subpoint", value))
        index = _ROMANS.index(value)
        if index + 1 < len(_ROMANS):
            next_subpoint = _ROMANS[index + 1]
        else:
            next_subpoint = None
\end{lstlisting}

\subsection{Text that closes a list}

A paragraph that introduces a list of points often continues after the list
with text that applies to the whole paragraph, not to its last point. EU
drafting calls it a \emph{subparagraph}. Article 101(1) of the AI Act lists
four grounds for a fine in points (a) to (d) and then states, without a marker
of its own, how the amount of the fine is fixed. Nothing in the text marks
where point (d) ends, so a parser that only reads markers would give the
fining criteria the label of point (d)---and a question about point (d) would
receive them too.

The end of a list is recognised from the punctuation of formal drafting. Items
before the last end with a semicolon, a comma or a conjunction; only the last
ends a sentence. So after the last item of a list, a full stop followed by a
line that begins with a capital letter starts the closing text. Inside a
numbered paragraph the closing text is labelled as the paragraph's second
subparagraph (third, and so on, after further lists); outside a numbered
paragraph there is no subparagraph to name, and the closing text simply
returns to the level of the provision. This also covers documents other than
legislation: without it, the guidance that follows the conditions listed in
TCPS~2 Article 3.7A would carry the label of its last condition.
\Cref{lst:article101} shows the labels the chunks of Article 101 receive.

\begin{lstlisting}[float, numbers=none, caption={Labels assigned to the chunks of
Article 101 of the AI Act.}, label={lst:article101}]
p.117  Article 101               Article 101 Fines for providers of ...
p.117  Article 101(1)            1. The Commission may impose ... fines
p.117  Article 101(1)(a)         (a) infringed the relevant provisions ...
p.117  Article 101(1)(b)         (b) failed to comply with a request ...
p.117  Article 101(1)(c)         (c) failed to comply with a measure ...
p.118  Article 101(1)(d)         (d) failed to make available ...
p.118  Article 101(1), 2nd sub.  In fixing the amount of the fine ...
p.118  Article 101(2)            2. Before adopting the decision ...
\end{lstlisting}

\paragraph{Trade-offs.} The rule relies on drafting conventions rather than
on layout, so it works on text extracted from any source. It would close a
list too early if its last item ran over several sentences with a line break
after the first, which formal drafting rarely does. Articles whose lists are
not followed by closing text, such as Article 7 of the AI Act, are labelled
exactly as before.

\subsection{Asking for a subdivision}

The question is parsed for the subdivision it names (\cref{tab:subdivisions}).
A single letter is a point and two or more roman letters are a sub-point, so
``point (i)'' reads as a point and ``point (iv)'' as a sub-point. A number
after \emph{point} or \emph{item} is read as a paragraph, because that is how
people refer to numbered paragraphs in practice.

\begin{table}[ht]
\centering
\small
\begin{tabular}{@{}ll@{}}
\toprule
Written in the question & Read as \\
\midrule
paragraph 2, para 2, subsection 2, 66(2)  & paragraph 2 \\
point d, letter d, (d)                    & point d \\
point (iv), subpoint iv, (iv)             & sub-point iv \\
point 3, item 3                           & paragraph 3 \\
second subparagraph, subparagraph 2       & subparagraph 2 \\
\bottomrule
\end{tabular}
\caption{Subdivisions recognised in a question.}
\label{tab:subdivisions}
\end{table}

Exact lookup (\cref{sec:retrieval}) fetches the whole provision, and the
chunks that carry every requested label are then kept; a paragraph's first
subparagraph is its text without a subparagraph label. If no chunk carries
them all, the outer labels are dropped one at a time, so ``66(2)(e)'' still
finds point (e) in an article whose points are not grouped into numbered
paragraphs. Only the matching chunks are passed on: ``paragraph 2 of Article
7'' returns 14 chunks---the paragraph's opening, points (a) to (k), and the two
sub-points of point (k)---while ``point (k) of paragraph 2 of Article 7''
returns three, and ``the second subparagraph of paragraph 1 of Article 101''
returns the closing text alone. When the provision has no such subdivision, the
lookup note lists what it has instead, which produces the deterministic reply
described in \cref{sec:generation}.

\subsection{Writing references for the reader}

A passage's labels are written back the way the reference would be cited:
the provision, then the paragraph, point and sub-point in parentheses, then the
subparagraph in words---``Article 66(2)(e)(iv)'', ``Article 101(1), second
subparagraph''---preceded by the printed page where the document has one. The
same wording is used for the focus line given to the model
(\cref{sec:generation}), so the model and the reader see the same reference.
